%\section{Former equational theories}\label{appendix:oldgraphicallanguage}

\begin{figure}[ht]
    % \scalebox{.73}{\doublebox{\begin{minipage}{1.338\textwidth}
    \fbox{\begin{minipage}{.975\textwidth}\begin{center}
        \vspace{-1em}
        \begin{subfigure}{0.21\textwidth}
            \begin{align}\label{C1}\tag{C1}\tikzfig{cliffordaxioms/w}=\tikzfig{gates/empty}\end{align}
        \end{subfigure}
        \begin{subfigure}{0.29\textwidth}
            \begin{align}\label{C2}\tag{C2}\tikzfig{cliffordaxioms/HH}=\tikzfig{gates/Id}\end{align}
        \end{subfigure}
        \begin{subfigure}{0.37\textwidth}
            \begin{align}\label{C3}\tag{C3}\tikzfig{cliffordaxioms/SSSS}=\tikzfig{gates/Id}\end{align}
        \end{subfigure}\vspace{-0.5em}
        \begin{subfigure}{0.52\textwidth}
            \begin{align}\label{C4}\tag{C4}\tikzfig{cliffordczaxioms/C4-00}=\tikzfig{cliffordczaxioms/C4-03}\end{align}
        \end{subfigure}
        \begin{subfigure}{0.28\textwidth}
            \begin{align}\label{C5}\tag{C5}\tikzfig{cliffordczaxioms/C5-00}=\tikzfig{identities/I2}\end{align}
        \end{subfigure}\vspace{-0.5em}
        \begin{subfigure}{0.32\textwidth}
            \begin{align}\label{C6}\tag{C6}\tikzfig{cliffordczaxioms/C6-00}=\tikzfig{cliffordczaxioms/C6-03}\end{align}
        \end{subfigure}
        \begin{subfigure}{0.32\textwidth}
            \begin{align}\label{C7}\tag{C7}\tikzfig{cliffordczaxioms/C7-00}=\tikzfig{cliffordczaxioms/C7-05}\end{align}
        \end{subfigure}\vspace{-0.5em}
        \begin{subfigure}{0.61\textwidth}
            \begin{align}\label{C8}\tag{C8}\tikzfig{cliffordczaxioms/C8-00}=\tikzfig{cliffordczaxioms/C8-08}\end{align}
        \end{subfigure}\vspace{-0.5em}
        \begin{subfigure}{0.61\textwidth}
            \begin{align}\label{C9}\tag{C9}\tikzfig{cliffordczaxioms/C9-00}=\tikzfig{cliffordczaxioms/C9-09}\end{align}
        \end{subfigure}\vspace{-0.5em}
        \begin{subfigure}{0.61\textwidth}
            \begin{align}\label{C10}\tag{C10}\tikzfig{cliffordczaxioms/C10-00}=\tikzfig{cliffordczaxioms/C10-14}\end{align}
        \end{subfigure}\vspace{-0.5em}
        \begin{subfigure}{0.61\textwidth}
            \begin{align}\label{C11}\tag{C11}\tikzfig{cliffordczaxioms/C11-00}=\tikzfig{cliffordczaxioms/C11-01}\end{align}
        \end{subfigure}
        \begin{subfigure}{0.35\textwidth}
            \begin{align}\label{C12}\tag{C12}\tikzfig{cliffordczaxioms/C12-00}=\tikzfig{cliffordczaxioms/C12-06}\end{align}
        \end{subfigure}
        \begin{subfigure}{.85\textwidth}
            \begin{align}\label{C13}\tag{C13}\tikzfig{cliffordczaxioms/C13-00}=\tikzfig{cliffordczaxioms/C13-25}\end{align}
        \end{subfigure}
        \begin{subfigure}{.82\textwidth}
            \begin{align}\label{C14}\tag{C14}\tikzfig{cliffordczaxioms/C14-00}=\tikzfig{cliffordczaxioms/C14-09}\end{align}
        \end{subfigure}
        \begin{subfigure}{.82\textwidth}
            \begin{align}\label{C15}\tag{C15}\tikzfig{cliffordczaxioms/C15-00}=\tikzfig{cliffordczaxioms/C15-01}\end{align}
        \end{subfigure}
        \begin{subfigure}{.44\textwidth}
            \begin{align}\label{DSWAPdef}\tag{S}\tikzfig{gates/SWAP}=\tikzfig{cliffordczaxioms/DSWAPdef-02}\end{align}
        \end{subfigure}
    \end{center}\end{minipage}}
    % \end{minipage}}}
    \caption{Set of relations $\QCcliffordCZ$. Equations $(\textup{C}_i)$ are the ones of the complete set of relations for the pro considered in \cite{SelingerStabilizer}.}
    \label{appendix:oldgraphicallanguage1}
\end{figure}

\begin{figure}[ht]
    \fbox{\begin{minipage}{.975\textwidth}\begin{center}
        \vspace{-1em}
        \begin{subfigure}{0.23\textwidth}
            \begin{align}\label{R1}\tag{R1}\tikzfig{realcliffordaxioms/2old} = \tikzfig{gates/empty}\end{align}
        \end{subfigure}
        \begin{subfigure}{0.28\textwidth}
            \begin{align}\label{R2}\tag{R2}\tikzfig{realcliffordaxioms/Z2} = \tikzfig{gates/Id}\end{align}
        \end{subfigure}
        \begin{subfigure}{0.28\textwidth}
            \begin{align}\label{R3}\tag{R3}\tikzfig{realcliffordaxioms/H2} = \tikzfig{gates/Id}\end{align}
        \end{subfigure}\vspace{-0.5em}
        \begin{subfigure}{0.63\textwidth}
            \begin{align}\label{R4}\tag{R4}\tikzfig{realcliffordczaxioms/R4-00}=\tikzfig{realcliffordczaxioms/R4-02}\end{align}
        \end{subfigure}
        \begin{subfigure}{0.27\textwidth}
            \begin{align}\label{R5}\tag{R5}\tikzfig{cliffordczaxioms/C5-00}=\tikzfig{identities/I2}\end{align}
        \end{subfigure}\vspace{-0.5em}
        \begin{subfigure}{0.34\textwidth}
            \begin{align}\label{R6}\tag{R6}\tikzfig{realcliffordczaxioms/R6-00}=\tikzfig{realcliffordczaxioms/R6-07}\end{align}
        \end{subfigure}
        \begin{subfigure}{0.34\textwidth}
            \begin{align}\label{R7}\tag{R7}\tikzfig{realcliffordczaxioms/R7-00}=\tikzfig{realcliffordczaxioms/R7-01}\end{align}
        \end{subfigure}\vspace{-0.5em}
        \begin{subfigure}{0.52\textwidth}
            \begin{align}\label{R8}\tag{R8}\tikzfig{realcliffordczaxioms/R8-00}=\tikzfig{realcliffordczaxioms/R8-06}\end{align}
        \end{subfigure}\vspace{-0.5em}
        \begin{subfigure}{0.52\textwidth}
            \begin{align}\label{R9}\tag{R9}\tikzfig{realcliffordczaxioms/R9-00}=\tikzfig{realcliffordczaxioms/R9-01}\end{align}
        \end{subfigure}\vspace{-0.5em}
        \begin{subfigure}{0.5\textwidth}
            \begin{align}\label{R10}\tag{R10}\tikzfig{realcliffordczaxioms/R10-00}=\tikzfig{realcliffordczaxioms/R10-10}\end{align}
        \end{subfigure}\vspace{-0.5em}
        \begin{subfigure}{0.5\textwidth}
            \begin{align}\label{R11}\tag{R11}\tikzfig{realcliffordczaxioms/R11-00}=\tikzfig{realcliffordczaxioms/R11-01}\end{align}
        \end{subfigure}\vspace{-0.5em}
        \begin{subfigure}{0.64\textwidth}
            \begin{align}\label{R12}\tag{R12}\tikzfig{realcliffordczaxioms/R12-00}=\tikzfig{realcliffordczaxioms/R12-12}\end{align}
        \end{subfigure}
        \begin{subfigure}{0.34\textwidth}
            \begin{align}\label{R13}\tag{R13}\tikzfig{realcliffordczaxioms/R13-00}=\tikzfig{realcliffordczaxioms/R13-99}\end{align}
        \end{subfigure}\vspace{-0.5em}
        \begin{subfigure}{1\textwidth}
            \begin{align}\label{R14}\tag{R14}\tikzfig{realcliffordczaxioms/R14-00}=\tikzfig{realcliffordczaxioms/R14-99}\end{align}
        \end{subfigure}\vspace{-0.5em}
        \begin{subfigure}{.76\textwidth}
            \begin{align}\label{R15}\tag{R15}\tikzfig{realcliffordczaxioms/R15-00}=\tikzfig{realcliffordczaxioms/R15-99}\end{align}
        \end{subfigure}\vspace{-0.5em}
        \begin{subfigure}{.76\textwidth}
            \begin{align}\label{R16}\tag{R16}\tikzfig{realcliffordczaxioms/R16-00}=\tikzfig{realcliffordczaxioms/R16-99}\end{align}
        \end{subfigure}\vspace{-0.5em}
        \begin{subfigure}{.46\textwidth}
            \begin{align}\label{real-DSWAPdef}\tag{S}\tikzfig{gates/SWAP}=\tikzfig{cliffordczaxioms/DSWAPdef-04}\end{align}
        \end{subfigure}
    \end{center}\end{minipage}}
    \caption{Set of relations $\QCrealcliffordCZ$. Equations $(\textup{R}_i)$ are the ones of the complete set of relations for the pro considered in \cite{SelingerRealStabilizer}.}
    \label{appendix:oldgraphicallanguage2}
\end{figure}



\begin{figure}[ht]
    \fbox{\begin{minipage}{.975\textwidth}\begin{center}
        \vspace{-1em}
        \begin{subfigure}{0.26\textwidth}
            \begin{align}\label{qt-C1}\tag{C1}
                \tikzfig{qutritcliffordaxioms/omegapow6}
                =
                \tikzfig{gates/empty}
            \end{align}
        \end{subfigure}
        %
        \begin{subfigure}{0.36\textwidth}
            \begin{align}\label{qt-C2}\tag{C2}
                \tikzfig{qutritcliffordaxioms/hadpow4}
                =
                \tikzfig{gates/Id}
            \end{align}
        \end{subfigure}
        %
        \begin{subfigure}{0.32\textwidth}
            \begin{align}\label{qt-C3}\tag{C3}
                \tikzfig{qutritcliffordaxioms/spow3}
                =
                \tikzfig{gates/Id}
            \end{align}
        \end{subfigure}\vspace{-0.5em}
        %
        \begin{subfigure}{0.48\textwidth}
            \begin{align}\label{qt-C4}\tag{C4}
                \tikzfig{qutritcliffordaxioms/s2hs2hs2h}
                =
                \tikzfig{qutritcliffordaxioms/Iphase-omega}
            \end{align}
        \end{subfigure}
        %
        \begin{subfigure}{0.49\textwidth}
            \begin{align}\label{qt-C5}\tag{C5}
                \tikzfig{qutritcliffordaxioms/ssprime}
                =
                \tikzfig{qutritcliffordaxioms/sprimes}
            \end{align}
        \end{subfigure}\vspace{-0.5em}
        %
        \begin{subfigure}{0.36\textwidth}
            \begin{align}\label{qt-C6}\tag{C6}
                \tikzfig{qutritderivations/CZ-CZ-CZ} = \tikzfig{identities/I2}
            \end{align}
        \end{subfigure}
        %
        \begin{subfigure}{0.44\textwidth}
            \begin{align}\label{qt-SWAPSWAP}\tag{C7}
                \tikzfig{qutritderivations/SWAP-SWAP} = \tikzfig{identities/I2}
            \end{align}
        \end{subfigure}\vspace{-0.5em}
        \begin{subfigure}{0.36\textwidth}
            \begin{align}\label{qt-C7}\tag{C8}
                \tikzfig{qutritderivations/S-CZ} = \tikzfig{qutritderivations/CZ-S}
            \end{align}
        \end{subfigure}
        %
        \begin{subfigure}{0.41\textwidth}
            \begin{align}\label{qt-C8}\tag{C9}
                \tikzfig{qutritderivations/K-CZ}
                =
                \tikzfig{qutritderivations/CZ-CZ-K}
            \end{align}
        \end{subfigure}\vspace{-0.5em}
        %
        \begin{subfigure}{0.47\textwidth}
            \begin{align}\label{qt-SSWAP}\tag{C10}
                \tikzfig{qutritderivations/S-SWAP} = \tikzfig{qutritderivations/SWAP-S}
            \end{align}
        \end{subfigure}
        \begin{subfigure}{0.47\textwidth}
            \begin{align}\label{qt-C9}\tag{C11}
                \tikzfig{qutritderivations/H-SWAP} = \tikzfig{qutritderivations/SWAP-H}
            \end{align}
        \end{subfigure}\vspace{-0.5em}
        %
        \begin{subfigure}{0.47\textwidth}
            \begin{align}\label{qt-CZSWAP}\tag{C12}
                \tikzfig{qutritderivations/CZ-SWAP} = \tikzfig{qutritderivations/SWAP-CZz}
            \end{align}
        \end{subfigure}
        %
        \begin{subfigure}{0.5\textwidth}
            \begin{align}\label{qt-C10}\tag{C13}
                \tikzfig{qutritderivations/S-NOTC-origin} = \tikzfig{qutritderivations/NOTC-CZ-SSp}
            \end{align}
        \end{subfigure}
        %
        \begin{subfigure}{0.54\textwidth}
            \begin{align}\label{qt-C11}\tag{C14}
                \tikzfig{qutritderivations/CZ-CNOT-b} = \tikzfig{qutritderivations/NOTC-CZ-SSp-b}
            \end{align}
        \end{subfigure}
        %
        \begin{subfigure}{0.38\textwidth}
            \begin{align}\label{qt-C13}\tag{C15}
                \tikzfig{qutritderivations/ICZ-CZI} = \tikzfig{qutritderivations/CZI-ICZ}
            \end{align}
        \end{subfigure}
        %
        \begin{subfigure}{0.67\textwidth}
            \begin{align}\label{qt-C14}\tag{C16}
                \tikzfig{qutritderivations/ISW-SWI-ISW} = \tikzfig{qutritderivations/SWI-ISW-SWI}
            \end{align}
        \end{subfigure}
        %
        \begin{subfigure}{0.63\textwidth}
            \begin{align}\label{qt-C15}\tag{C17}
                \tikzfig{qutritderivations/ICZ-SWI-ISW} = \tikzfig{qutritderivations/SWI-ISW-CZI}
            \end{align}
        \end{subfigure}
        %
        \begin{subfigure}{0.41\textwidth}
            \begin{align}\label{qt-C16}\tag{C18}
                \tikzfig{qutritderivations/CNI-ICZ} = \tikzfig{qutritderivations/ICZ-CNI-CIZ}
            \end{align}
        \end{subfigure}
        %
        \begin{subfigure}{0.38\textwidth}
            \begin{align}\label{qt-Swap}\tag{S}
                \tikzfig{qutritderivations/SWAP-CZ} = \tikzfig{gates/SWAP}
            \end{align}
        \end{subfigure}
        
    \end{center}\end{minipage}}
    \caption{Set of relations $\QCqutritcliffordCZ$. Equations $(\textup{C}_i)$ are the ones of the complete set of relations for the pro considered in \cite{qutritclifford}.}
    \label{appendix:oldgraphicallanguage3}
\end{figure}



\begin{figure}[ht]
    % \scalebox{.73}{\doublebox{\begin{minipage}{1.338\textwidth}
    \fbox{\begin{minipage}{.975\textwidth}\begin{center}
        \vspace{-1em}
        \begin{subfigure}{0.21\textwidth}
            \begin{align}\label{t-C3}\tag{C3}\tikzfig{cliffordaxioms/w}=\tikzfig{gates/empty}\end{align}
        \end{subfigure}
        \begin{subfigure}{0.29\textwidth}
            \begin{align}\label{t-C4}\tag{C4}\tikzfig{cliffordaxioms/HH}=\tikzfig{gates/Id}\end{align}
        \end{subfigure}
        \begin{subfigure}{0.37\textwidth}
            \begin{align}\label{t-C5}\tag{C5}\tikzfig{cliffordaxioms/SSSS}=\tikzfig{gates/Id}\end{align}
        \end{subfigure}\vspace{-0.5em}
        \begin{subfigure}{0.52\textwidth}
            \begin{align}\label{t-C6}\tag{C6}\tikzfig{cliffordczaxioms/C4-00}=\tikzfig{cliffordczaxioms/C4-03}\end{align}
        \end{subfigure}
        \begin{subfigure}{0.27\textwidth}
            \begin{align}\label{t-C14}\tag{C14}\tikzfig{shortcut/S}=\tikzfig{shortcut/TT}\end{align}
        \end{subfigure}\vspace{-0.5em}
        \begin{subfigure}{0.43\textwidth}
            \begin{align}\label{t-C15}\tag{C15}\tikzfig{cliffordplustaxioms/TX} = \tikzfig{cliffordplustaxioms/XTd}\end{align}
        \end{subfigure}
        \begin{subfigure}{0.28\textwidth}
            \begin{align}\label{t-C7}\tag{C7}\tikzfig{cliffordczaxioms/C5-00}=\tikzfig{identities/I2}\end{align}
        \end{subfigure}\vspace{-0.5em}
        \begin{subfigure}{0.32\textwidth}
            \begin{align}\label{t-C8}\tag{C8}\tikzfig{cliffordczaxioms/C6-00}=\tikzfig{cliffordczaxioms/C6-03}\end{align}
        \end{subfigure}
        \begin{subfigure}{0.32\textwidth}
            \begin{align}\label{t-C9}\tag{C9}\tikzfig{cliffordczaxioms/C7-00}=\tikzfig{cliffordczaxioms/C7-05}\end{align}
        \end{subfigure}\vspace{-0.5em}
        \begin{subfigure}{0.62\textwidth}
            \begin{align}\label{t-C10}\tag{C10}\tikzfig{cliffordczaxioms/C8-00}=\tikzfig{cliffordczaxioms/C8-08}\end{align}
        \end{subfigure}
        \begin{subfigure}{0.63\textwidth}
            \begin{align}\label{t-C11}\tag{C11}\tikzfig{cliffordczaxioms/C9-00}=\tikzfig{cliffordczaxioms/C9-09}\end{align}
        \end{subfigure}
        \begin{subfigure}{0.62\textwidth}
            \begin{align}\label{t-C12}\tag{C12}\tikzfig{cliffordczaxioms/C10-00}=\tikzfig{cliffordczaxioms/C10-14}\end{align}
        \end{subfigure}
        \begin{subfigure}{0.62\textwidth}
            \begin{align}\label{t-C13}\tag{C13}\tikzfig{cliffordczaxioms/C11-00}=\tikzfig{cliffordczaxioms/C11-01}\end{align}
        \end{subfigure}
        \begin{subfigure}{0.39\textwidth}
            \begin{align}\label{t-C16}\tag{C16}\tikzfig{cliffordplustaxioms/C16l}=\tikzfig{cliffordplustaxioms/C16r}\end{align}
        \end{subfigure}
        \begin{subfigure}{0.69\textwidth}
            \begin{align}\label{t-C17}\tag{C17}\tikzfig{cliffordplustaxioms/C17l}=\tikzfig{cliffordplustaxioms/C17r}\end{align}
        \end{subfigure}
        \begin{subfigure}{0.84\textwidth}
            \begin{align}\label{t-C18}\tag{C18}\tikzfig{cliffordplustaxioms/C18l}=\tikzfig{cliffordplustaxioms/C18r}\end{align}
        \end{subfigure}
        \begin{subfigure}{0.95\textwidth}
            \begin{align}\label{old-t-C19}\tag{C19}\scalebox{0.8}{\tikzfig{cliffordplustaxioms/C19l}}=\scalebox{0.8}{\tikzfig{cliffordplustaxioms/C19r}}\end{align}
        \end{subfigure}
        \begin{subfigure}{0.69\textwidth}
            \begin{align}\label{old-t-C20}\tag{C20}\tikzfig{cliffordplustaxioms/C20l}=\tikzfig{cliffordplustaxioms/C20r}\end{align}
        \end{subfigure}
        \begin{subfigure}{.44\textwidth}
            \begin{align}\label{t-DSWAPdef}\tag{S}\tikzfig{gates/SWAP}=\tikzfig{cliffordczaxioms/DSWAPdef-02}\end{align}
        \end{subfigure}
    \end{center}\end{minipage}}
    % \end{minipage}}}
    \caption{Set of relations $\QCcliffordplustCZ$ on two qubits. Equations $(\textup{C}_i)$ are the ones of the complete set of relations for the pro considered in \cite{SelingerCliffordPlusT}.}
    \label{appendix:oldgraphicallanguage4}
\end{figure}




\begin{figure}[ht]
    \scalebox{1}{
    \fbox{\begin{minipage}{.975\textwidth}\begin{center}
        \vspace{-1em}
        \begin{subfigure}{0.2\textwidth}
            \begin{align}\label{cs-C1}\tag{C1}\tikzfig{qutritcliffordaxioms/ipow4} = \tikzfig{gates/empty}\end{align}
        \end{subfigure}
        \begin{subfigure}{0.34\textwidth}
            \begin{align}\label{cs-C2}\tag{C2}\tikzfig{cliffordpluscsaxioms/K2}=\tikzfig{cliffordpluscsaxioms/ipow3}\end{align}
        \end{subfigure}
        \begin{subfigure}{0.4\textwidth}
            \begin{align}\label{cs-C3}\tag{C3}\tikzfig{cliffordaxioms/SSSS}=\tikzfig{cliffordaxioms/I}\end{align}
        \end{subfigure}
        \begin{subfigure}{0.52\textwidth}
            \begin{align}\label{cs-C4}\tag{C4}\tikzfig{cliffordpluscsaxioms/SKSKSK}=\tikzfig{cliffordpluscsaxioms/ipow3}\end{align}
        \end{subfigure}
        \begin{subfigure}{0.38\textwidth}
            \begin{align}\label{cs-C5}\tag{C5}\tikzfig{cliffordpluscsaxioms/CSpow4}=\tikzfig{identities/I2}\end{align}
        \end{subfigure}
        \begin{subfigure}{0.35\textwidth}
            \begin{align}\label{cs-C6}\tag{C6}\tikzfig{cliffordpluscsaxioms/ScommTopLeft}=\tikzfig{cliffordpluscsaxioms/ScommTopRight}\end{align}
        \end{subfigure}
        \begin{subfigure}{0.38\textwidth}
            \begin{align}\label{cs-C7}\tag{C7}\tikzfig{cliffordpluscsaxioms/ScommBotLeft}=\tikzfig{cliffordpluscsaxioms/ScommBotRight}\end{align}
        \end{subfigure}
        \begin{subfigure}{0.57\textwidth}
            \begin{align}\label{cs-C8}\tag{C8}\tikzfig{cliffordpluscsaxioms/XTopLeft}=\tikzfig{cliffordpluscsaxioms/XTopRight}\end{align}
        \end{subfigure}
        \begin{subfigure}{0.57\textwidth}
            \begin{align}\label{cs-C9}\tag{C9}\tikzfig{cliffordpluscsaxioms/XBotLeft}=\tikzfig{cliffordpluscsaxioms/XBotRight}\end{align}
        \end{subfigure}
        \begin{subfigure}{0.59\textwidth}
            \begin{align}\label{cs-C10}\tag{C10}\tikzfig{cliffordpluscsaxioms/KTopLeft}=\tikzfig{cliffordpluscsaxioms/KTopRight}\end{align}
        \end{subfigure}
        \begin{subfigure}{0.59\textwidth}
            \begin{align}\label{cs-C11}\tag{C11}\tikzfig{cliffordpluscsaxioms/KBotLeft}=\tikzfig{cliffordpluscsaxioms/KBotRight}\end{align}
        \end{subfigure}
        \begin{subfigure}{0.37\textwidth}
            \begin{align}\label{cs-C12}\tag{C12}\tikzfig{cliffordpluscsaxioms/ICSCSI}=\tikzfig{cliffordpluscsaxioms/CSIICS}\end{align}
        \end{subfigure}
        \begin{subfigure}{0.59\textwidth}
            \begin{align}\label{cs-C13}\tag{C13}\tikzfig{cliffordpluscsaxioms/C13lred}=\tikzfig{cliffordpluscsaxioms/C13rred}\end{align}
        \end{subfigure}
        \begin{subfigure}{0.67\textwidth}
            \begin{align}\label{cs-C14-old}\tag{C14}\tikzfig{cliffordpluscsaxioms/C14lred}=\tikzfig{cliffordpluscsaxioms/C14rred}\end{align}
        \end{subfigure}
        \begin{subfigure}{0.68\textwidth}
            \begin{align}\label{cs-C15}\tag{C15}\tikzfig{cliffordpluscsaxioms/C15lred}=\tikzfig{cliffordpluscsaxioms/C15rred}\end{align}
        \end{subfigure}
        \begin{subfigure}{0.75\textwidth}
            \begin{align}\label{cs-C16-old}\tag{C16}\tikzfig{cliffordpluscsaxioms/C16l}=\tikzfig{cliffordpluscsaxioms/C16r}\end{align}
        \end{subfigure}
        \begin{subfigure}{0.92\textwidth}
            \begin{align}\label{cs-C17-old}\tag{C17}\tikzfig{cliffordpluscsaxioms/C17l}=\tikzfig{cliffordpluscsaxioms/C17r}\end{align}
        \end{subfigure}
        \begin{subfigure}{0.62\textwidth}
            \begin{align}\tag{B}\label{cs-DSWAP}\tikzfig{gates/SWAP}=\tikzfig{cliffordpluscsczaxioms/DSWAP-final}\end{align}
        \end{subfigure}
    \end{center}\end{minipage}}
    }
    \caption{Set of relations $\QCcliffordpluscsCZ$ on three qubits. Equations $(\textup{C}_i)$ are the ones of the complete set of relations for the pro considered in \cite{SelingerCliffordPlusCS}.}
    \label{appendix:oldgraphicallanguage5}
\end{figure}











\begin{figure}[ht]
    \centering
    \fbox{%
        \begin{minipage}{0.98\textwidth}
        \centering
        % \begin{minipage}[c]{0.34\textwidth}
        %     \centering
        %     \begin{equation}\tag{TX}\label{cnot-R11}
        %     \scalebox{0.9}{\tikzfig{cnotdihedral/R11L}}
        %     =
        %     \scalebox{0.9}{\tikzfig{cnotdihedral/R11R}}
        %     \end{equation}
        % \end{minipage}
        % \begin{minipage}[c]{0.34\textwidth}
        %     \centering
        %     \begin{equation}\tag{XC}\label{new-R3}
        %     \scalebox{0.9}{\tikzfig{identities/wCNOT-00}}
        %     =
        %     \scalebox{0.9}{\tikzfig{identities/wCNOT-16}}
        %     \end{equation}
        % \end{minipage}
        % \begin{minipage}[c]{0.36\textwidth}
        %     \centering
        %     \begin{equation}\tag{ZC}\label{new-R8}
        %     \scalebox{0.9}{\tikzfig{identities/zCNOT-00}}
        %     =
        %     \scalebox{0.9}{\tikzfig{identities/zCNOT-05}}
        %     \end{equation}
        % \end{minipage}
        \begin{minipage}[c]{0.31\textwidth}
            \centering
            \begin{equation}\tag{R1}
            \scalebox{0.9}{\tikzfig{cnotdihedral/R1L}}
            =
            \scalebox{0.9}{\tikzfig{gates/Id}}
            \end{equation}
        \end{minipage}
        \begin{minipage}[c]{0.36\textwidth}
            \centering
            \begin{equation}\tag{R2}\label{cnot-old-R2}
            \scalebox{0.9}{\tikzfig{identities/CXC_00}}
            =
            \scalebox{0.9}{\tikzfig{identities/CXC_06}}
            \end{equation}
        \end{minipage}
        \begin{minipage}[c]{0.31\textwidth}
            \centering
            \begin{equation}\tag{R3}\label{cnot-old-R3}
            \scalebox{0.9}{\tikzfig{cliffordczaxioms/CNOTXCNOT}}
            =
            \scalebox{0.9}{\tikzfig{cliffordczaxioms/XX}}
            \end{equation}
        \end{minipage}
        \begin{minipage}[c]{0.29\textwidth}
            \centering
            \begin{equation}\tag{R4}\label{cnot-old-R4}
            \scalebox{0.9}{\tikzfig{identities/CNOTCNOT}}
            =
            \scalebox{0.9}{\tikzfig{identities/I2}}
            \end{equation}
        \end{minipage}
        \begin{minipage}[c]{0.47\textwidth}
            \centering
            \begin{equation}\tag{R5}\label{cnot-old-R5}
            \scalebox{0.9}{\tikzfig{cliffordczaxioms/CNOTswCNOTswCNOT}}
            =
            \scalebox{0.9}{\tikzfig{gates/SWAP}}
            \end{equation}
        \end{minipage}
        \begin{minipage}[c]{0.36\textwidth}
            \centering
            \begin{equation}\tag{R6}
            \scalebox{0.9}{\tikzfig{cnotdihedral/R6L}}
            =
            \scalebox{0.9}{\tikzfig{cnotdihedral/R6R}}
            \end{equation}
        \end{minipage}
        \begin{minipage}[c]{0.28\textwidth}
            \centering
            \begin{equation}\tag{R7}
            \scalebox{0.9}{\tikzfig{cnotdihedral/R7L}}
            =
            \scalebox{0.9}{\tikzfig{gates/Id}}
            \end{equation}
        \end{minipage}
        \begin{minipage}[c]{0.36\textwidth}
            \centering
            \begin{equation}\tag{R8}\label{cnot-old-R8}
            \scalebox{0.9}{\tikzfig{cliffordczaxioms/R8-old-decomp}}
            =
            \scalebox{0.9}{\tikzfig{cliffordczaxioms/ZZ}}
            \end{equation}
        \end{minipage}
        \begin{minipage}[c]{0.53\textwidth}
            \centering
            \begin{equation}\tag{R9}
            \scalebox{0.9}{\tikzfig{cnotdihedral/R9L}}
            =
            \scalebox{0.9}{\tikzfig{cnotdihedral/R9R}}
            \end{equation}
        \end{minipage}
        \begin{minipage}[c]{0.22\textwidth}
            \centering
            \begin{equation}\tag{R10}
                \scalebox{0.9}{\tikzfig{gates/wc}}^{\otimes 8} = \scalebox{0.9}{\tikzfig{gates/empty}}
            \end{equation}
        \end{minipage}
        \begin{minipage}[c]{0.37\textwidth}
            \centering
            \begin{equation}\tag{R11}\label{cnot-old-R11}
            \scalebox{0.9}{\tikzfig{cnotdihedral/R11oldL}}
            =
            \scalebox{0.9}{\tikzfig{cnotdihedral/R11oldR}}
            \end{equation}
        \end{minipage}
        \begin{minipage}[c]{0.33\textwidth}
            \centering
            \begin{equation}\tag{R12}
            \scalebox{0.9}{\tikzfig{cliffordplustaxioms/CNOTTCNOT}}
            =
            \scalebox{0.9}{\tikzfig{cliffordplustaxioms/T}}
            \end{equation}
        \end{minipage}
        \begin{minipage}[c]{0.9\textwidth}
            \centering
            \begin{equation}\tag{R13}
            \scalebox{0.82}{\tikzfig{cnotdihedral/R13L}}
            =
            \scalebox{0.82}{\tikzfig{cnotdihedral/R13R}}
            \end{equation}
        \end{minipage}
        \end{minipage}%
    }
    \caption{Set of relations $\QColddihedral$ \cite{AmyCNOTDihedral}.}
    \label{appendix:oldgraphicallanguage6}
\end{figure}





\begin{table}[t]
\centering
\small
\renewcommand{\arraystretch}{1.15}
%\setlength{\tabcolsep}{9pt}

% Space before/after horizontal rules (adjust 0.7ex to taste)
\newcommand{\spacedhline}{\noalign{\vskip 0.7ex}\hline\noalign{\vskip 0.7ex}}

\begin{tabular}{@{}l p{0.35\linewidth} p{0.35\linewidth}@{}}
\hline
\noalign{\vskip 0.7ex}

\textbf{Fragment} &
\textbf{Encoding $E_\bullet$ (new $\to$ src)} &
\textbf{Decoding $D_\bullet$ (src $\to$ new)} \\
\spacedhline

$\mathit{Cliff}, \mathit{RCliff}$ &
\(
\begin{aligned}
E_\bullet\left(\tikzfig{gates/CNOT}\right) &= \tikzfig{completeness/H2CZH2}
\end{aligned}
\) &
\(
\begin{aligned}
D_\bullet\left(\tikzfig{gates/CZ}\right) &= \tikzfig{completeness/H2CNOTH2}
\end{aligned}
\) \\
\spacedhline

$\mathit{CliffT}$ &
\(
\begin{aligned}
E_{\mathit{CliffT}}\left(\tikzfig{gates/CNOT}\right) &= \tikzfig{completeness/H2CZH2}
\end{aligned}
\) &
\(
\begin{aligned}
D_{\mathit{CliffT}}\left(\tikzfig{gates/CZ}\right) &= \tikzfig{completeness/H2CNOTH2}\\
D_{\mathit{CliffT}}\left(\tikzfig{gates/S}\right) &= \tikzfig{shortcut/TT}
\end{aligned}
\) \\
\spacedhline

$\mathit{Cliff3}$ &
\(
\begin{aligned}
E_{\mathit{Cliff3}}\left(\tikzfig{gates/S}\right) &= \tikzfig{gates/Sbis}\\
E_{\mathit{Cliff3}}\left(\tikzfig{gates/H}\right) &= \tikzfig{gates/Hbisp}\\
E_{\mathit{Cliff3}}\left(\tikzfig{gates/CNOT}\right) &= \tikzfig{completeness/H2CZHHH2}\\
E_{\mathit{Cliff3}}\left(\tikzfig{gates/wb}\right) &= \tikzfig{gates/wb}
\end{aligned}
\) &
\(
\begin{aligned}
D_{\mathit{Cliff3}}\left(\tikzfig{gates/S}\right) &= \tikzfig{gates/Sbis}\\
D_{\mathit{Cliff3}}\left(\tikzfig{gates/H}\right) &= \tikzfig{gates/Hbis}\\
D_{\mathit{Cliff3}}\left(\tikzfig{gates/CZ}\right) &= \tikzfig{completeness/HHH2CNOTH2}\\
D_{\mathit{Cliff3}}\left(\tikzfig{gates/wb}\right) &= \tikzfig{gates/wb}\\
D_{\mathit{Cliff3}}\left(\tikzfig{gates/minw}\right) &= \tikzfig{gates/w10}
\end{aligned}
\) \\
\spacedhline

$\mathit{CliffCS}$ &
\(
\begin{aligned}
E_{\mathit{CliffCS}}\left(\tikzfig{gates/H}\right) &= \tikzfig{gates/Kw}\\
E_{\mathit{CliffCS}}\left(\tikzfig{gates/w}\right) &= \tikzfig{gates/w}
\end{aligned}
\) &
\(
\begin{aligned}
D_{\mathit{CliffCS}}\left(\tikzfig{gates/K}\right) &= \tikzfig{gates/Hw}\\
D_{\mathit{CliffCS}}\left(\tikzfig{gates/i}\right) &= \tikzfig{gates/w2}\\
D_{\mathit{CliffCS}}\left(\tikzfig{gates/w}\right) &= \tikzfig{gates/w}
\end{aligned}
\) \\
\noalign{\vskip 0.7ex}
\hline

\end{tabular}

\caption{Encoding/decoding pairs used in the completeness-transfer arguments.
For the qutrit Clifford and Clifford+CS rows, the source presentation is the
scalar-refined one from \cref{lem:scalar-refinement}.
\emph{Convention:} omitted generators are fixed by $E_\bullet$ and $D_\bullet$.}
\label{tab:ED-clifford-family}
\end{table}

